\subsection{Conclusión individual: Emanuel Pino}

Mi principal aporte en esta práctica fue preparar el plan de intercambios, ejecutar las corridas de TCP y gRPC y verificar los eventos registrados en los archivos CSV. Al revisar cómo se habían ejecutado las primeras mediciones, encontré una diferencia entre la topología utilizada y la que pedía la consigna. Por eso repetí la comparativa con tres procesos independientes en localhost y conservé los datos anteriores. Registrar el PID y el puerto de cada nodo me permitió comprobar cómo se ejecutó la prueba y dejar esa información disponible para quien quiera reproducirla.

La verificación de Lamport confirmó que, en los 9\,600 pares de envío y recepción de la nueva comparativa, la recepción siempre tuvo una marca mayor que el envío. También encontré dos solicitudes de nodos distintos que salieron con la misma marca, \(C=1957\). Ese ejemplo me ayudó a entender que la marca de Lamport respeta la precedencia causal, pero puede coincidir entre eventos concurrentes.

Al revisar las tablas, observé que TCP tuvo menor latencia media en las cuatro condiciones medidas. Con 100 bytes y tres emisores concurrentes, el RTT medio fue de 0,416\,ms en TCP y de 0,680\,ms en gRPC. Los resultados corresponden a una sola máquina. Si repitiera la práctica, registraría la topología y las versiones desde la primera corrida y probaría también en máquinas separadas para observar cómo cambia la comparación cuando interviene una red real.
